\daytitle{27}{近代物理与原子物理入门}

\begin{learningbox}{本讲义大致内容}
近代物理从经典物理解释不了的现象出发：光电效应说明光具有粒子性；原子能级说明原子内部能量不是连续变化；原子核反应说明质量和能量可以相互联系。本节目标不是展开大学量子力学，而是建立“微观世界常常是不连续的、概率性的、守恒约束非常强”的基本认识。
\end{learningbox}

\section*{一、学习目标}
\begin{enumerate}
  \item 理解光电效应的实验事实和爱因斯坦光电方程。
  \item 掌握光子能量 $E=h\nu$、逸出功和截止频率的含义。
  \item 理解原子能级跃迁与光谱的关系。
  \item 会配平核反应方程，掌握质量数和电荷数守恒。
  \item 理解半衰期和质能方程的基本意义。
\end{enumerate}

\section*{二、初高中衔接}

前面的光学主要把光看成光线或波。近代物理告诉我们，光在某些现象中必须看成一份一份的光子。前面的能量守恒仍然成立，但微观系统的能量变化常常不是连续的，而是按确定能级差吸收或放出。

\textbf{科学发展视角：} 光电效应是太阳能电池、光电传感器和数码相机感光元件的基础；能级跃迁支撑激光、光谱分析和天体成分识别；核物理连接核能、医学影像、放射性测年和粒子物理。

\section*{三、核心概念}

\begin{conceptbox}{光电效应}

用一定频率的光照射金属表面，金属中电子吸收光子能量后可能逸出，形成光电子。实验表明：能否发生光电效应主要由光的频率决定，而不是由光强单独决定。

\[
\boxed{h\nu=W_0+E_k}
\]

这条式子可以看成微观世界的能量账本：一个光子的能量先用于把电子从金属中“拉出来”，这部分能量叫逸出功；如果还有剩余，才表现为光电子的最大初动能。因此频率不够时，即使光很强，也只能增加照来的光子数，不能让单个光子具有更大的能量。

\begin{center}
% TikZ 图：光电效应
% Day 27：近代物理
\begin{tikzpicture}[scale=0.9,>=Latex]
  \draw[fill=gray!20,thick] (0,0) rectangle (3.8,0.35);
  \node at (1.9,-0.3) {金属板};
  \foreach \x in {0.6,1.25,1.9}{
    \draw[->,blue,thick] (\x,2.25) -- (\x+0.4,0.45);
    \node[blue] at (\x-0.1,2.45) {$h\nu$};
  }
  \foreach \x/\ang in {1.0/55,1.7/75,2.4/42}{
    \draw[->,red,thick] (\x,0.36) -- ++({1.0*cos(\ang)},{1.0*sin(\ang)});
    \fill[red] (\x,0.36) circle (1.6pt);
  }
  \node[red] at (3.75,1.45) {光电子};
  \node[align=left,font=\small] at (6.25,1.1)
    {$h\nu=W_0+E_{k\max}$\\[4pt]
     $\nu\ge\nu_0$ 时\\能发生光电效应。};
\end{tikzpicture}


\figcaption{光电效应：一个光子的能量被一个电子吸收}\label{fig:photoelectric-effect}
\end{center}

\begin{itemize}
  \item $h\nu$：光子能量。
  \item $W_0$：金属逸出功。
  \item $E_k$：逸出光电子的最大初动能。
  \item 截止频率满足 $h\nu_0=W_0$。
\end{itemize}
\end{conceptbox}

\begin{conceptbox}{原子能级}

原子只能处在一系列特定能量状态中，这些能量状态叫能级。原子吸收或放出光子时，能量变化满足
\[
\boxed{h\nu=|E_m-E_n|}.
\]

\begin{center}
% TikZ 图：原子能级跃迁
% Day 27：近代物理
\begin{tikzpicture}[scale=0.9,>=Latex]
  \foreach \y/\lab in {0/$E_1$,1.1/$E_2$,1.85/$E_3$,2.35/$E_4$}{
    \draw[thick] (0,\y) -- (4.45,\y);
    \node[left] at (0,\y) {\lab};
  }
  \draw[->,red,thick] (1.0,2.35) -- (1.0,1.1);
  \draw[->,blue,thick] (2.1,0) -- (2.1,1.85);
  \draw[->,red,thick] (3.1,1.85) -- (3.1,0);
  \node[red,anchor=west] at (4.65,1.55) {放出光子};
  \node[blue,anchor=west] at (4.65,0.75) {吸收光子};
  \draw[red!70,thin] (3.18,1.05) -- (4.52,1.55);
  \draw[blue!70,thin] (2.18,1.42) -- (4.52,0.75);
  \node[font=\small] at (2.25,-0.55) {$h\nu=|E_m-E_n|$};
\end{tikzpicture}


\figcaption{原子能级跃迁：能级差决定吸收或放出光子的频率}\label{fig:energy-level-transition}
\end{center}

\textbf{理解：} 能级不是电子绕核运动的普通轨道线，而是原子允许具有的能量状态。能级差不同，发出的光频率不同，于是形成特征光谱。
\end{conceptbox}

\begin{conceptbox}{原子核与核反应}

原子核由质子和中子组成。核反应方程必须满足两个守恒：
\[
\boxed{\text{质量数守恒}},\qquad
\boxed{\text{电荷数守恒}}.
\]

\begin{center}
% TikZ 图：核反应方程守恒
% Day 27：近代物理
\begin{tikzpicture}[scale=0.9,>=Latex]
  \node[draw,rounded corners=2pt,fill=blue!8,minimum width=2.3cm,minimum height=1.0cm] (left) at (0,0) {反应前};
  \node[draw,rounded corners=2pt,fill=red!8,minimum width=2.3cm,minimum height=1.0cm] (right) at (4.1,0) {反应后};
  \draw[->,thick] (left) -- (right);
  \node[align=left,font=\small] at (8.0,0.3)
    {质量数守恒：$\sum A_{\text{前}}=\sum A_{\text{后}}$\\[4pt]
     电荷数守恒：$\sum Z_{\text{前}}=\sum Z_{\text{后}}$};
  \node[font=\small] at (4.10,-0.85)
    {$Q=(m_{\text{前}}-m_{\text{后}})c^2$：$Q>0$ 释放，$Q<0$ 吸收};
\end{tikzpicture}


\figcaption{核反应方程：配平的核心是质量数和电荷数守恒}\label{fig:nuclear-reaction-balance}
\end{center}

常见符号：
\[
{}^4_2\mathrm{He}\quad(\alpha\text{ 粒子}),\qquad
{}^0_{-1}e\quad(\beta^-\text{ 粒子}),\qquad
{}^1_0n\quad(\text{中子}).
\]
\end{conceptbox}

\begin{conceptbox}{半衰期}

大量放射性原子核中，有一半发生衰变所需的时间叫半衰期。半衰期是统计规律，不是说某一个原子核到时间一定衰变。

\begin{center}
% TikZ 图：半衰期曲线
% Day 27：近代物理
\begin{tikzpicture}[scale=0.9,>=Latex]
  \draw[->,thick] (0,0) -- (5.4,0) node[right] {$t$};
  \draw[->,thick] (0,0) -- (0,3.3) node[above] {$N$};
  \draw[blue,thick,samples=80,domain=0:5] plot(\x,{2.8*exp(-0.55*\x)});
  \foreach \x/\lab/\y in {0/$0$/2.8,1.26/$T_{1/2}$/1.4,2.52/$2T_{1/2}$/0.7,3.78/$3T_{1/2}$/0.35}{
    \draw[dashed,gray] (\x,0) -- (\x,\y) -- (0,\y);
    \node[below] at (\x,0) {\lab};
  }
  \node[left] at (0,2.8) {$N_0$};
  \node[left] at (0,1.4) {$N_0/2$};
  \node[left] at (0,0.7) {$N_0/4$};
  \node[font=\scriptsize,align=center,blue!70!black] at (3.95,2.55) {每过一个\\半衰期减半};
\end{tikzpicture}


\figcaption{半衰期：剩余原子核数按固定比例减少}\label{fig:half-life-curve}
\end{center}
\end{conceptbox}

\section*{四、公式与适用条件}

\begin{center}
\begin{tabular}{lll}
\toprule
公式 & 物理意义 & 适用条件 \\
\midrule
$E=h\nu$ & 光子能量 & 单个光子 \\
$h\nu=W_0+E_k$ & 光电方程 & 最大初动能 \\
$h\nu_0=W_0$ & 截止频率 & 刚好逸出 \\
$h\nu=|E_m-E_n|$ & 能级跃迁 & 原子吸收或放光 \\
$E=mc^2$ & 质能关系 & 静能量 \\
$\Delta E=\Delta mc^2$ & 核能释放 & 质量亏损 \\
\bottomrule
\end{tabular}
\end{center}

\section*{五、典型模型}

\subsection*{模型 1：光电效应判断}

先比较频率：
\[
\nu<\nu_0\Rightarrow\text{不能发生光电效应},
\]
\[
\nu>\nu_0\Rightarrow\text{能发生光电效应}.
\]
光强增大只会增加单位时间逸出电子数，不会改变单个光子的能量。

\subsection*{模型 2：能级跃迁}

向高能级跃迁需要吸收光子；向低能级跃迁会放出光子。光子频率由能级差决定：
\[
\nu=\frac{|E_m-E_n|}{h}.
\]

\subsection*{模型 3：核反应配平}

未知粒子可通过质量数、电荷数分别列方程判断。例如
\[
{}^{14}_{7}\mathrm{N}+{}^4_2\mathrm{He}\rightarrow{}^{17}_{8}\mathrm{O}+X.
\]
质量数：$14+4=17+A_X$，得 $A_X=1$；电荷数：$7+2=8+Z_X$，得 $Z_X=1$，所以
\[
X={}^1_1\mathrm{H}.
\]

\section*{六、例题精讲}

\begin{examplebox}{光电效应最大初动能}
\textbf{题意分析：} 某金属逸出功 $W_0=2.0\ \mathrm{eV}$，用光子能量为 $3.5\ \mathrm{eV}$ 的光照射，求光电子最大初动能。

\textbf{物理模型：} 爱因斯坦光电方程。

\begin{center}
% TikZ 图：光电效应例题的能量核算
% Day 27：近代物理
\begin{tikzpicture}[scale=0.86,>=Latex,line cap=round,line join=round]
  \path[use as bounding box] (-0.25,-0.90) rectangle (7.70,1.55);
  \draw[->,blue,thick] (0.05,0.35) -- (1.45,0.35);
  \node[blue,font=\small,fill=white,inner sep=1pt] at (0.80,0.82)
    {$h\nu=3.5\ \mathrm{eV}$};
  \draw[thick,fill=gray!16] (1.70,-0.05) rectangle (2.50,0.75);
  \node[font=\small] at (2.10,0.35) {金属};
  \draw[->,red,thick] (2.75,0.35) -- (4.25,0.35);
  \node[red,font=\small,fill=white,inner sep=1pt] at (3.50,0.82)
    {$W_0=2.0\ \mathrm{eV}$};
  \draw[->,orange,thick] (4.55,0.35) -- (6.45,0.35);
  \node[orange,font=\small,fill=white,inner sep=1pt] at (5.52,0.82)
    {$E_{k\max}$};
  \node[font=\small,align=center] at (3.90,-0.55)
    {光子能量先克服逸出功，剩余部分成为最大初动能};
\end{tikzpicture}

\end{center}

\textbf{解题过程：}
\[
E_k=h\nu-W_0=3.5-2.0=1.5\ \mathrm{eV}.
\]

\textbf{方法总结：} 光电效应题先看频率或光子能量是否超过逸出功，再算最大初动能。
\end{examplebox}

\begin{examplebox}{能级跃迁发光频率}
\textbf{题意分析：} 原子从 $E_3=-1.5\ \mathrm{eV}$ 跃迁到 $E_1=-13.6\ \mathrm{eV}$，求放出光子的能量。

\textbf{物理模型：} 能级差等于光子能量。

\begin{center}
% TikZ 图：能级跃迁例题
% Day 27：近代物理
\begin{tikzpicture}[scale=0.88,>=Latex,line cap=round,line join=round]
  \draw[thick] (0,0) -- (3.80,0);
  \node[left] at (0,0) {$E_1=-13.6\ \mathrm{eV}$};
  \draw[thick] (0,2.00) -- (3.80,2.00);
  \node[left] at (0,2.00) {$E_3=-1.5\ \mathrm{eV}$};
  \draw[->,red,thick] (2.15,2.00) -- (2.15,0.08);
  \draw[->,orange,thick] (2.35,1.35) -- (4.40,1.35)
    node[right] {$E_\gamma=12.1\ \mathrm{eV}$};
  \node[font=\small,align=center] at (2.00,-0.48)
    {高能级降到低能级，能级差以光子形式放出};
\end{tikzpicture}

\end{center}

\textbf{解题过程：}
\[
E_{\gamma}=E_3-E_1=(-1.5)-(-13.6)=12.1\ \mathrm{eV}.
\]

\textbf{方法总结：} 能级为负值时，要用高能级减低能级；放出光子能量为正。
\end{examplebox}

\begin{examplebox}{半衰期计算}
\textbf{题意分析：} 某放射性样品初始有 $N_0$ 个未衰变原子核，半衰期为 5 天。经过 15 天后剩余多少？

\textbf{物理模型：} 每过一个半衰期剩余一半。

\begin{center}
% TikZ 图：半衰期阶梯计算
% Day 27：近代物理
\begin{tikzpicture}[scale=0.88,>=Latex,line cap=round,line join=round]
  \path[use as bounding box] (-0.45,-0.65) rectangle (6.90,1.55);
  \foreach \x/\lab/\val in {0/0/$N_0$,1.8/5/$\frac{N_0}{2}$,3.6/10/$\frac{N_0}{4}$,5.4/15/$\frac{N_0}{8}$}{
    \draw[fill=blue!8,thick] (\x,0) rectangle ++(1.05,0.70);
    \node at (\x+0.525,0.35) {\val};
    \node[font=\small] at (\x+0.525,-0.32) {\lab 天};
  }
  \foreach \x in {1.18,2.98,4.78}{
    \draw[->,thick] (\x,0.36) -- ++(0.42,0);
    \node[font=\scriptsize,above] at (\x+0.21,0.42) {$\times\frac12$};
  }
\end{tikzpicture}

\end{center}

\textbf{解题过程：}
15 天包含
\[
\frac{15}{5}=3
\]
个半衰期，所以
\[
N=N_0\left(\frac12\right)^3=\frac{N_0}{8}.
\]

\textbf{方法总结：} 半衰期是比例衰减，不是每隔一个半衰期减少同样数量。
\end{examplebox}

\section*{七、变式训练}

\begin{enumerate}
  \item 某金属截止频率为 $\nu_0$，若入射光频率为 $0.8\nu_0$，能否发生光电效应？
  \item 光强增大时，光电子最大初动能是否一定增大？说明理由。
  \item 原子从高能级跃迁到低能级时，是吸收还是放出光子？
  \item 配平核反应：${}^{238}_{92}\mathrm{U}\rightarrow{}^{234}_{90}\mathrm{Th}+X$，判断 $X$。
\end{enumerate}

\subsection*{训练提示}
\begin{enumerate}
  \item 不能。入射光频率低于截止频率，单个光子能量不足。
  \item 不一定。若频率不变，最大初动能不变；光强增大主要使光电子数增多。
  \item 从高能级到低能级放出光子；从低能级到高能级吸收光子。
  \item 质量数差为 4，电荷数差为 2，所以 $X={}^{4}_{2}\mathrm{He}$。
\end{enumerate}

\section*{八、易错点}
\begin{mistakebox}
\begin{enumerate}
  \item \textbf{认为光强越大一定能打出电子：} 是否能发生光电效应取决于频率是否超过截止频率。
  \item \textbf{混淆光强和光子能量：} 光强增大通常意味着光子数增多，不代表单个光子能量增大。
  \item \textbf{能级差符号处理错误：} 光子能量取能级差的绝对值。
  \item \textbf{核反应只配质量数：} 必须同时配平质量数和电荷数。
  \item \textbf{把半衰期用于单个原子核：} 半衰期是大量原子核的统计规律。
\end{enumerate}
\end{mistakebox}

\section*{九、竞赛入门补充}
\begin{competitionbox}
\textbf{德布罗意物质波}

近代物理的一个重要思想是波粒二象性：光有粒子性，微观粒子也有波动性。德布罗意提出，动量为 $p$ 的粒子对应波长
\[
\lambda=\frac{h}{p}.
\]

\textbf{为什么宏观物体看不出波动性？}
宏观物体动量很大，所以对应波长极小，远小于可观察尺度。电子等微观粒子动量小，波长可与原子尺度相当，因此会表现出衍射和干涉。

\begin{center}
% TikZ 图：德布罗意物质波与电子衍射
% Day 27：近代物理
\begin{tikzpicture}[scale=0.88,>=Latex,line cap=round,line join=round]
  \path[use as bounding box] (-0.35,-1.05) rectangle (7.10,1.65);
  \draw[->,blue,thick] (0,0.35) -- (1.55,0.35)
    node[midway,above] {电子束};
  \draw[thick] (2.00,-0.65) -- (2.00,1.35);
  \foreach \y in {-0.35,0.35,1.05}{
    \draw[fill=white] (2.00,\y) circle (0.08);
  }
  \node[font=\small] at (2.00,-0.92) {晶体};
  \draw[red,thick] (2.15,0.35) .. controls (3.10,1.10) and (4.00,1.18) .. (5.15,0.92);
  \draw[red,thick] (2.15,0.35) .. controls (3.10,-0.40) and (4.00,-0.48) .. (5.15,-0.22);
  \draw[red,thick] (2.15,0.35) -- (5.15,0.35);
  \draw[thick] (5.45,-0.75) -- (5.45,1.45);
  \foreach \y in {-0.22,0.35,0.92}{
    \fill[red] (5.45,\y) circle (1.8pt);
  }
  \node[font=\small,align=center] at (6.20,0.35) {衍射\\亮斑};
  \node[font=\small] at (3.50,-0.85) {$\lambda=\dfrac{h}{p}$};
\end{tikzpicture}

\end{center}

\textbf{连接前后知识：} 这一公式把动量 $p$、普朗克常量 $h$ 和波长 $\lambda$ 联系起来，是从经典物理走向量子物理的桥梁。
\end{competitionbox}

\section*{十、模型深化：实验图像与现代应用}

\subsection*{1.\quad 光电效应实验装置}

光电效应不是“看见电子飞出来”这么简单，而是通过光电流、电压和光照条件的关系来判断微观过程。实验中用光照射阴极，若有光电子逸出并到达阳极，外电路中就出现光电流。

\begin{center}
% TikZ 图：光电效应实验装置
% Day 27：近代物理
\begin{tikzpicture}[scale=0.9,>=Latex]
  \draw[thick,rounded corners] (-3,-1.4) rectangle (3,1.4);
  \node[above] at (0,1.4) {抽成真空的光电管};
  \draw[thick] (-1.8,-0.8) -- (-1.8,0.8);
  \node[below] at (-1.8,-0.82) {阴极 K};
  \draw[thick] (1.8,-0.8) -- (1.8,0.8);
  \node[right] at (1.8,0) {阳极 A};
  \foreach \y in {-0.45,0,0.45} {
    \draw[->,blue,thick] (-3.55,\y+0.35) -- (-1.95,\y);
  }
  \node[blue] at (-3.75,1.05) {入射光};
  \draw[->,red,thick] (-1.6,0.2) -- (1.45,0.45) node[midway,above] {光电子};
  \draw[->,red,thick] (-1.6,-0.25) -- (1.45,-0.15);
  \draw[thick] (-1.8,-1.4) -- (-1.8,-2.0) -- (-0.95,-2.0);
  \draw[thick] (-0.70,-2.0) circle (0.25);
  \node at (-0.70,-2.0) {$A$};
  \draw[thick] (-0.45,-2.0) -- (0.72,-2.0);
  \draw[thick] (0.72,-2.30) -- (0.72,-1.70);
  \draw[thick] (1.02,-2.16) -- (1.02,-1.84);
  \draw[thick] (1.02,-2.0) -- (1.8,-2.0) -- (1.8,-1.4);
  \node[font=\scriptsize] at (-1.35,-2.45) {微安表};
  \node[font=\scriptsize] at (1.45,-2.45) {可调电源（可反接）};
  \node[font=\small,align=center] at (0,-2.92)
    {调节外加电压，观察光电流是否存在以及何时为零};
\end{tikzpicture}


\figcaption{光电效应实验装置：用光电流和外加电压研究光电子能量}\label{fig:photoelectric-experiment}
\end{center}

\textbf{实验事实到方程：} 截止电压说明最快光电子的初动能可被反向电场刚好“耗尽”，所以
\[
eU_c=E_{k\max}.
\]
改变入射光频率，发现 $E_{k\max}$ 随频率线性增大，由此得到
\[
h\nu=W_0+E_{k\max}.
\]

\subsection*{2.\quad 光电效应伏安曲线}

光电效应实验不仅能说明“有无光电子”，还可以通过伏安曲线研究光电子最大初动能和单位时间逸出电子数。

\begin{center}
% TikZ 图：光电效应伏安特性
% Day 27：近代物理
\begin{tikzpicture}[scale=0.9,>=Latex]
  \draw[->,thick] (-2.4,0) -- (4.0,0) node[right] {$U$};
  \draw[->,thick] (0,0) -- (0,3.0) node[above] {$I$};
  \coordinate (Uc) at (-1.35,0);
  \draw[blue,thick] (Uc) .. controls (-0.72,1.08) and (0.45,2.15) .. (3.3,2.25);
  \draw[red,thick] (Uc) .. controls (-0.55,0.72) and (0.65,1.48) .. (3.3,1.62);
  \draw[dashed] (Uc) -- (-1.35,-0.34);
  \node[below] at (-1.35,-0.34) {$-U_c$};
  \node[blue] at (3.55,2.25) {强光};
  \node[red] at (3.55,1.65) {弱光};
  \node[font=\small,align=center] at (2.05,2.75) {同频率、不同光强};
  \node[font=\small,align=center] at (1.65,-1.05) {同频率时截止电压相同；饱和电流反映光电子数};
\end{tikzpicture}


\figcaption{光电效应伏安曲线：截止电压和饱和电流对应不同物理意义}\label{fig:photoelectric-ui-graph}
\end{center}

\begin{itemize}
  \item \textbf{截止电压：} 使光电流刚好为 0 的反向电压，满足 $eU_c=E_{k\max}$。
  \item \textbf{饱和电流：} 所有逸出的光电子都被收集时的电流，反映单位时间逸出电子数。
  \item \textbf{光强影响：} 同频率下，光强越大，饱和电流越大；但截止电压不因光强增大而增大。
\end{itemize}

\subsection*{3.\quad 核能释放的数量级}

核反应中的质量亏损通常很小，但由于
\[
c^2\approx9.0\times10^{16}\ \mathrm{m^2/s^2},
\]
即使很小的质量亏损也可能对应巨大的能量。这解释了为什么核裂变和核聚变能释放远高于普通化学反应的能量。

\textbf{严谨提醒：} 质量亏损不是“质量消失了所以不守恒”的简单说法。更准确地说，在相对论中质量和能量共同参与守恒，核反应前后的总能量守恒。

\begin{examplebox}{截止电压与最大初动能}
\textbf{题意分析：} 某光电管实验中，截止电压为 $1.2\ \mathrm{V}$。求光电子最大初动能（用 eV 表示）。

\textbf{物理模型：} 反向电场做功恰好阻止最快光电子到达阳极。

\textbf{解题过程：}
\[
E_{k\max}=eU_c.
\]
若用电子伏特作单位，电荷量为 $e$ 的电子通过 $1\ \mathrm{V}$ 电势差对应 $1\ \mathrm{eV}$，所以
\[
E_{k\max}=1.2\ \mathrm{eV}.
\]

\textbf{方法总结：} 用 eV 处理光电效应能大大简化单位换算。
\end{examplebox}

\begin{examplebox}{核反应方程补全}
\textbf{题意分析：} 补全核反应：
\[
{}^{27}_{13}\mathrm{Al}+{}^4_2\mathrm{He}\rightarrow{}^{30}_{15}\mathrm{P}+X.
\]

\textbf{物理模型：} 质量数守恒、电荷数守恒。

\textbf{解题过程：}
质量数：
\[
27+4=30+A_X,\qquad A_X=1.
\]
电荷数：
\[
13+2=15+Z_X,\qquad Z_X=0.
\]
所以
\[
X={}^1_0n.
\]

\textbf{方法总结：} 未知粒子先由 $A$、$Z$ 判断，再识别名称。
\end{examplebox}

\subsection*{4.\quad 分层训练}

\begin{enumerate}
  \item \textbf{基础：} 入射光频率增大，光电子最大初动能怎样变化？
  \item \textbf{提高：} 同频率光强增大，饱和光电流和截止电压分别怎样变化？
  \item \textbf{综合：} 某原子从 $E_4$ 跃迁到 $E_2$ 放出光子，再从 $E_2$ 跃迁到 $E_1$ 放出光子，比较两光子能量应看什么？
  \item \textbf{开放思考：} 为什么核聚变是太阳能量来源，却很难在地面稳定受控实现？
\end{enumerate}

\section*{十一、今日小结}
\begin{enumerate}
  \item 光电效应说明光具有粒子性，光子能量 $E=h\nu$。
  \item 光电方程为 $h\nu=W_0+E_k$，截止频率满足 $h\nu_0=W_0$。
  \item 原子能级跃迁满足 $h\nu=|E_m-E_n|$。
  \item 核反应方程必须满足质量数守恒和电荷数守恒。
  \item 半衰期是大量原子核衰变的统计规律。
\end{enumerate}
